\section{Common envelopes and dominated continuity}

\label{sec:relative-common-envelope-setting}
Let $u_j\in\mathcal E_\varphi$ converge in $L^1$, with measures dominated by one finite non-pluripolar measure $\nu$. To keep the limit in $\mathcal E_\varphi$, we construct, after passing to a subsequence, a common lower envelope in this class. Partitioning the contact set of a finite rooftop according to the obstacle it touches shows that the rooftop measure is again dominated by $\nu$, with no factor depending on the number of obstacles. A summable integration error, combined with the lower mass bound $\kappa$, then prevents the stretched rooftops from collapsing, uniformly in the number of obstacles.

\begin{lemma}\label{lem:dominated-relative-rooftop-extraction}\index{Envelope!common lower envelope}
Suppose $u_j\in\mathcal E_\varphi$, $u_j\to u\in\PSH(X,\theta)$ in $L^1(X)$, and $\mu_{u_j}\leq\nu$ for a fixed finite non-pluripolar measure $\nu$. There is a subsequence, denoted again by $u_j$, such that
\[
 h=P_\theta\left(\inf_j u_j\right)\in\mathcal E_\varphi.
\]
In particular $u\in\mathcal E_\varphi$.
\end{lemma}
\begin{proof}
We first bound stretched rooftops of the tail of the sequence from below and then add the finitely many remaining terms.

The suprema of an $L^1$-convergent sequence are bounded, so \eqref{eq:relative-normalized-model-bound} gives a constant $C$ with $u_j\leq\varphi+C$ for all $j$. After a common translation we may assume $u_j\leq\varphi\leq0$, and then $u\leq\varphi$. Since $\mu_{u_j}\leq\nu$, \cref{lem:fixed-npp-integrals} gives
\[
 I_j=\int_X|\mathrm{e}^{u_j}-\mathrm{e}^u|\,\mathrm{d}\mu_{u_j}\longrightarrow0.
\]
We extract, once and for all, a subsequence with $I_j\leq2^{-j}$ along which $u_j\to u$ almost everywhere for smooth volume.

\Needspace{8\baselineskip}
For $N\geq j$ put $p_{j,N}=P_\theta(u_j,\ldots,u_N)$. These rooftops belong to $\mathcal E_\varphi$ by \cref{prop:relative-finite-rooftops}. Partition their contact sets according to the first index at which the minimum is attained. Ordered contact on these disjoint pieces gives both
\begin{equation}\label{eq:finite-rooftop-common-domination}
 \mu_{p_{j,N}}\leq\nu,
 \qquad
 \int_X|\mathrm{e}^{p_{j,N}}-\mathrm{e}^u|\,\mathrm{d}\mu_{p_{j,N}}
 \leq\sum_{\ell=j}^N I_\ell\leq2^{1-j}.
\end{equation}
The first bound has no factor $N$, because each piece is controlled by the restriction of the same measure $\nu$ to that piece.

Fix a real number $B\geq1$ and set $r_{j,N}=Q_B(p_{j,N},\varphi)$. It belongs to $\mathcal E_\varphi$ by the stretching lemma, so $\mu_{r_{j,N}}(X)\geq\kappa$. With $D_{j,N}=\{B^{-1}r_{j,N}+(1-B^{-1})\varphi=p_{j,N}\}$, contact support and positivity give
\[
 \mu_{r_{j,N}}\leq B^n\mathbf1_{D_{j,N}}\mu_{p_{j,N}}
 \leq B^n\nu.
\]
On this contact set, $\varphi\leq0$ gives
$\mathrm{e}^{p_{j,N}}\leq \mathrm{e}^{r_{j,N}/B}$. Consequently
\begin{equation}\label{eq:relative-rooftop-noncollapse-test}
 \int_X\mathrm{e}^u\,\mathrm{d}\mu_{r_{j,N}}
 \leq B^n2^{1-j}
      +B^n\nu(X)\exp\!\left(\frac{\sup_X r_{j,N}}B\right).
\end{equation}
On the other hand, choose $K>0$ with $B^n\nu(\{u\leq-K\})<\kappa/2$; this is possible because $u\not\equiv-\infty$ and $\nu$ does not charge its pole set. The measure domination then gives $\mu_{r_{j,N}}(\{u\leq-K\})<\kappa/2$, and the lower mass bound yields
\[
 \int_X\mathrm{e}^u\,\mathrm{d}\mu_{r_{j,N}}\geq \mathrm{e}^{-K}\kappa/2.
\]
Choose $j$ so large that $B^n2^{1-j}<\mathrm{e}^{-K}\kappa/4$. Then \eqref{eq:relative-rooftop-noncollapse-test} bounds $\sup_Xr_{j,N}$ from below, independently of $N$. Since the sequence decreases in $N$ and is bounded above by $\varphi$, its limit $R$ is not identically $-\infty$. It satisfies, for every $\ell\geq j$,
\begin{equation}\label{eq:relative-tail-stretched-competitor}
 B^{-1}R+(1-B^{-1})\varphi\leq u_\ell,
 \qquad R\leq\varphi.
\end{equation}
The second inequality follows already from $r_{j,N}\leq p_{j,N}\leq\varphi$.

It remains to include the finitely many earlier terms and to verify every defining condition of the relative class. Fix $b\geq1$, apply the preceding construction with $B=2b$, and let $j$ and $R$ be as above. If $j>1$, let $p_0=P_\theta(u_1,\ldots,u_{j-1})$ and $S=Q_{2b}(p_0,\varphi)$. By the finite-rooftop and stretching lemmas, neither is identically $-\infty$, and $S\leq\varphi$. If $j=1$, set $S=\varphi$. The potential $H=(R+S)/2$ is $\theta$-plurisubharmonic. For a tail index $\ell\geq j$ we use \eqref{eq:relative-tail-stretched-competitor} and $S\leq\varphi$; for an earlier index we use the analogous inequality for $S$ and $R\leq\varphi$. In either case,
\[
 b^{-1}H+(1-b^{-1})\varphi\leq u_\ell
 \quad\hbox{for every }\ell\geq1.
\]
For $b=1$ this shows that $h=P_\theta(\inf_\ell u_\ell)$ is not identically $-\infty$. For general $b$ the left side is a competitor for $h$, so the same inequality holds with $h$ in place of every $u_\ell$; this is the defining condition of the relative class for $b$. Also $h\leq\varphi$. Thus $h\in\mathcal E_\varphi$.

Finally $h\leq u_j$ for all $j$, so the almost-everywhere convergence and \cref{lem:qpsh-almost-everywhere-order} give $h\leq u\leq\varphi$. Upward stability proves $u\in\mathcal E_\varphi$. Undoing the common translation proves the lemma as stated.
\end{proof}

Once a common lower envelope is available, convergence in capacity follows from a single integral condition. The proof uses the zero-domination lemma, so infinite masses are allowed.

\Needspace{17\baselineskip}
\begin{proposition}\label{prop:relative-common-envelope-criterion}
Let $h\in\mathcal E_\varphi$, let $C$ be a constant, and let $u_j\in\PSH(X,\theta)$ satisfy
\[
 h\leq u_j\leq\varphi+C,\qquad
 u_j\longrightarrow u\in\PSH(X,\theta)\quad\hbox{in }L^1(X).
\]
If
\[
 I_j:=\int_X|\mathrm{e}^{u_j}-\mathrm{e}^u|\,\mu_{u_j}\longrightarrow0,
\]
then $u\in\mathcal E_\varphi$ and $u_j\to u$ in capacity.
\end{proposition}
\begin{proof}
The bounds pass to $u$ almost everywhere along a subsequence, and then everywhere by \cref{lem:qpsh-almost-everywhere-order}; upward stability gives $u_j,u\in\mathcal E_\varphi$. After translating the entire family and its bounds by a common constant, we may assume that all potentials are nonpositive.

It suffices to prove convergence in capacity along a further subsequence of every subsequence. Extract so that $I_j\leq2^{-j}$ and $u_j\to u$ almost everywhere for smooth volume, and form
\[
 p_{j,N}=P_\theta(u_j,\ldots,u_N),\qquad
 p_j=P_\theta\left(\inf_{\ell\geq j}u_\ell\right).
\]
All these envelopes lie above $h$ and below $\varphi+C$; hence none of them is identically $-\infty$, and all belong to $\mathcal E_\varphi$. Partition the contact set of $p_{j,N}$ according to the first index $\ell$ with $p_{j,N}=u_\ell$. Ordered contact on each part gives
\[
 \int_X|\mathrm{e}^{p_{j,N}}-\mathrm{e}^u|\,\mu_{p_{j,N}}
 \leq\sum_{\ell=j}^N I_\ell\leq2^{1-j}.
\]
Only integrals of nonnegative functions occur here, so the estimate holds even if a rooftop has infinite mass.

As $N\to\infty$, $p_{j,N}\downarrow p_j$; as $j\to\infty$, the $p_j$ increase, and their regularized limit $p$ is not identically $-\infty$. Both monotone limits hold in capacity by \cref{lem:monotone-psh-capacity}. Moreover $h\leq p\leq u\leq\varphi+C$, where $p\leq u$ follows from the almost-everywhere convergence of the extracted $u_j$. Thus $p\in\mathcal E_\varphi$.

Apply \cref{lem:weighted-extended-lsc}, first as $N\to\infty$ and then as $j\to\infty$, with the bounded quasi-continuous weights in the preceding estimate. These weights converge in capacity because the exponential is Lipschitz on a common upper half-line. We obtain
\[
 \int_X|\mathrm{e}^p-\mathrm{e}^u|\,\mu_p=0,
 \qquad \mu_p(\{p<u\})=0.
\]
The second equality follows off the poles, which the extended measures ignore. The zero-domination lemma \cref{lem:relative-zero-domination} now gives $p=u$; no finiteness of $\mu_p$ is needed.

Finally, the upper envelopes $U_j=(\sup_{\ell\geq j}u_\ell)^*$ decrease to $u$, by Hartogs' lemma (\cref{lem:qpsh-hartogs}). Applying \cref{lem:monotone-psh-capacity} to $p_j\uparrow u$ and $U_j\downarrow u$, the inequalities $p_j\leq u_j\leq U_j$ give convergence in capacity along this subsequence, and hence for the original sequence.
\end{proof}

\subsection{Convergence under a fixed measure}

The extraction lemma provides the lower envelope required by the integral criterion. The same fixed measure then controls the measures of deep sublevel sets, so convergence in capacity passes to the Monge--Amp\`ere measures themselves.

\begin{theorem}\label{thm:relative-fixed-measure-continuity}\label{thm:relative-dominated-continuity}\index{Non-pluripolar measure!dominated continuity}
Suppose $u_j\in\mathcal E_\varphi$, $u_j\to u\in\PSH(X,\theta)$ in $L^1(X)$, and $\mu_{u_j}\leq\nu$ for a fixed finite non-pluripolar measure $\nu$. Then
\[
 u\in\mathcal E_\varphi,
 \qquad u_j\longrightarrow u\text{ in capacity},
 \qquad \mu_{u_j}\rightharpoonup\mu_u\leq\nu.
\]
\end{theorem}
\begin{proof}
Start with an arbitrary subsequence. By \cref{lem:dominated-relative-rooftop-extraction}, a further subsequence has a common lower envelope $h\in\mathcal E_\varphi$, and $u\in\mathcal E_\varphi$. Inequality \eqref{eq:relative-normalized-model-bound} and Hartogs' lemma give a common upper bound $u_j\leq\varphi+C$.

After a common translation, assume $u_j,u\leq0$. Since $\mu_{u_j}\leq\nu$, \cref{lem:fixed-npp-integrals} gives
\[
 \int_X|\mathrm{e}^{u_j}-\mathrm{e}^u|\,\mathrm{d}\mu_{u_j}
 \leq\int_X|\mathrm{e}^{u_j}-\mathrm{e}^u|\,\mathrm{d}\nu
 \longrightarrow0.
\]
Thus \cref{prop:relative-common-envelope-criterion} gives convergence in capacity along this subsequence. For the measures themselves, the common lower bound gives
\[
 \mu_{u_j}(\{u_j\leq-T\})
 \leq\nu(\{h\leq-T\})\longrightarrow0
 \quad\hbox{uniformly in }j,
\]
since $\nu$ is finite and does not charge the pole set of $h\not\equiv-\infty$. The tail criterion for smooth backgrounds, \cref{thm:general-background-damping}, therefore gives $\mu_{u_j}\rightharpoonup\mu_u$, and weak limits preserve domination by the fixed finite measure $\nu$. Every subsequence of the original sequence has a further subsequence with these conclusions and the same limit $u$, so both convergences hold for the original sequence.
\end{proof}
